\scanpage{25}
\mspage{12}
d'examiner les actions de $\pi_1(f),\pi_1(g)$ sur les groupes abélianisés des $\pi_1$, et même sur leurs composantes $\ell$-adiques. Prenant une jacobienne généralisée du type ``ext. d'une V.A. par un tore'' de $X$, on sait que

1°) Les morphismes $f:X\longrightarrow G$ tels que $f(a)=0$ se factorisent de façon unique par
\diagram{\begin{tikzcd}[ampersand replacement=\&]X \arrow[r,"{\mathrm{can}}"] \& J \arrow[r,"\varphi"] \& G\end{tikzcd}}
avec $\varphi$ un hom. de groupes algébriques.

2°) \deleted{$\varphi$ est} Un tel hom. $\varphi$ est connu quand on connaît son action sur les $H_1(\ ,\mathbb Z_\ell)$. Ce qui revient à le connaître sur les points d'ordre $\ell^\nu$, quel que soit $\nu$ -- or ceux-ci sont denses\ldots

3°) $H_1(X,\mathbb Z_\ell)\simeq H_1(J,\mathbb Z_\ell)$

De ceci, \deleted{\unsure{le cas}} on conclut \inserted{(par 3°)} que $H_1(f)=H_1(g)$ implique (si $f=\varphi\circ\mathrm{can}$, $g=\psi\circ\mathrm{can}$) $H_1(\varphi)=H_1(\psi)$, donc par 2°) que $\varphi=\unsure{\varphi|\psi}$, donc $f=g$, cqfd.

\deleted{b)} Notons que l'on ne peut conclure de l'hyp. \inserted{de a)} que $f=g$, puisque pour $X=Y=G$, les translations $\tau_a$ induisent l'identité sur $\pi_1(G)$ (du moins \unsure{en la} car. $0$ ou si $G$ est une VA), sans être l'identité.

b) On va pourtant prouver l'égalité,
\margin{24}
